\section{Mo-PAE Residual-Risk Disparity}
\label{sec:residual-risk}

Table~\ref{tab:04d-decile-reduction}'s \emph{plain} and \emph{no-privacy}
rows report, for each dataset, the relative reduction in
re-identification risk (Definition~2) at the lowest- and
highest-entropy deciles, and the residual-risk-disparity ratio
(Definition~3) between them, under plain Mo-PAE's own privacy-utility
trade-off (no fairness term). Table~\ref{tab:inmodel-vs-external}
reports the mechanistic finding from the in-model-vs.-external
comparison described in Section~\ref{sec:methods-changed}.

\begin{table*}[t]
\centering
\caption{Checkpoint 04d: utility, relative risk reduction at deciles 0 and 9, and residual-risk disparity (RRD, post-mechanism risk decile 0 $\div$ decile 9), every approved-batch run. Bootstrap 95\% CIs in brackets; $\geq 3$ encoder seeds per point. Generated from \texttt{reports/figures/phase4d\_analysis/decile\_reduction\_metrics.json}.}
\label{tab:04d-decile-reduction}
\small
\begin{tabular}{@{}llcccc@{}}
\toprule
Dataset & Run & Utility & Rel.\ red., decile 0 & Rel.\ red., decile 9 & RRD (0/9) \\
\midrule
GeoLife & plain Mo-PAE & 0.532$\pm$0.052 & +0.014 [-0.133, +0.176] & +0.371 [-0.042, +0.625] & 8.08 [3.72, 14.11] \\
GeoLife & no-privacy ref. & 0.623$\pm$0.014 & -0.032 [-0.173, +0.104] & +0.340 [-0.045, +0.604] & 8.02 [3.79, 13.93] \\
GeoLife & ed\_soft $\lambda_4$=0.3 & 0.641$\pm$0.039 & +0.160 [+0.016, +0.295] & +0.239 [-0.076, +0.481] & 5.57 [2.67, 9.40] \\
GeoLife & ed\_soft $\lambda_4$=1.0 & 0.654$\pm$0.039 & +0.255 [+0.079, +0.408] & +0.245 [-0.151, +0.516] & 5.00 [2.49, 8.38] \\
GeoLife & ed\_soft $\lambda_4$=3.0 & 0.631$\pm$0.014 & +0.302 [+0.196, +0.419] & +0.299 [-0.019, +0.534] & 5.08 [2.44, 8.76] \\
GeoLife & floor $\lambda_4$=0.3 & 0.551$\pm$0.047 & -0.031 [-0.158, +0.096] & +0.347 [-0.097, +0.614] & 8.19 [3.95, 14.20] \\
GeoLife & floor $\lambda_4$=1.0 & 0.545$\pm$0.055 & -0.010 [-0.154, +0.114] & +0.425 [+0.142, +0.603] & 8.82 [4.56, 14.63] \\
GeoLife & floor $\lambda_4$=3.0 & 0.516$\pm$0.085 & -0.008 [-0.128, +0.111] & +0.473 [+0.210, +0.650] & 9.68 [4.74, 16.43] \\
\addlinespace
Foursquare & plain Mo-PAE & 0.111$\pm$0.005 & +0.384 [+0.331, +0.439] & +0.897 [+0.866, +0.925] & 16.02 [10.84, 23.71] \\
Foursquare & no-privacy ref. & 0.280$\pm$0.002 & +0.096 [+0.065, +0.127] & +0.310 [+0.250, +0.374] & 3.43 [2.74, 4.32] \\
Foursquare & ed\_soft $\lambda_4$=0.3 & 0.134$\pm$0.021 & +0.300 [+0.243, +0.358] & +0.850 [+0.810, +0.885] & 12.38 [8.61, 17.59] \\
Foursquare & ed\_soft $\lambda_4$=1.0 & 0.221$\pm$0.011 & +0.135 [+0.095, +0.179] & +0.600 [+0.543, +0.658] & 5.68 [4.37, 7.38] \\
Foursquare & ed\_soft $\lambda_4$=3.0 & 0.233$\pm$0.004 & +0.133 [+0.096, +0.171] & +0.568 [+0.508, +0.628] & 5.28 [4.10, 6.85] \\
\addlinespace
Gowalla-NYC & plain Mo-PAE & 0.269$\pm$0.035 & +0.240 [+0.115, +0.371] & +0.634 [+0.519, +0.731] & 7.26 [4.24, 11.71] \\
Gowalla-NYC & no-privacy ref. & 0.322$\pm$0.005 & +0.200 [+0.077, +0.343] & +0.334 [+0.245, +0.430] & 4.12 [2.63, 6.27] \\
Gowalla-NYC & ed\_soft $\lambda_4$=0.3 & 0.263$\pm$0.005 & +0.253 [+0.149, +0.360] & +0.555 [+0.452, +0.650] & 5.82 [3.61, 9.08] \\
Gowalla-NYC & ed\_soft $\lambda_4$=1.0 & 0.239$\pm$0.026 & +0.266 [+0.160, +0.381] & +0.596 [+0.527, +0.664] & 6.29 [3.86, 10.02] \\
Gowalla-NYC & ed\_soft $\lambda_4$=3.0 & 0.169$\pm$0.043 & +0.244 [+0.154, +0.345] & +0.629 [+0.540, +0.719] & 7.13 [4.24, 11.69] \\
\addlinespace
\bottomrule
\end{tabular}
\end{table*}

\begin{table}[t]
\centering
\caption{In-model adversary confidence vs.\ external attacker top-1, plain Mo-PAE, deciles 0 and 9. The in-model adversary stays far closer to chance and far flatter across deciles than the external attacker on the identical features. Generated from \texttt{reports/figures/phase4d\_analysis/inmodel\_vs\_external.json}.}
\label{tab:inmodel-vs-external}
\small
\begin{tabular}{@{}lccc@{}}
\toprule
Dataset & Chance & Decile 0 (ext.\ / in-model) & Decile 9 (ext.\ / in-model) \\
\midrule
GeoLife & 0.0055 & 0.345 / 0.0246 & 0.045 / 0.0113 \\
Foursquare & 0.0009 & 0.487 / 0.0521 & 0.031 / 0.0020 \\
Gowalla-NYC & 0.0001 & 0.566 / 0.0776 & 0.082 / 0.0063 \\
\bottomrule
\end{tabular}
\end{table}


\textbf{Plain Mo-PAE gives decile 0 no protection distinguishable from
none, on the dataset where this was checked most carefully.} On
GeoLife, decile 0's relative risk reduction under plain Mo-PAE is not
statistically distinguishable from zero (Table~\ref{tab:main-disparate-risk}),
while decile 9 receives a clearly positive reduction. The
residual-risk-disparity ratio is large: decile 0's post-mechanism
re-identification risk is several times decile 9's. On Foursquare and
Gowalla-NYC, plain Mo-PAE does give decile 0 a clearly nonzero
reduction ($+0.384$ and $+0.240$ respectively, both CIs excluding
zero), but the disparity is still large (residual-risk ratio
$16.02\times$ and $7.26\times$) because decile 9 is reduced far more
($+0.897$ and $+0.634$).

\textbf{A disparity-ratio claim and an absolute-difference claim can
disagree in direction, and this paper reports both for exactly that
reason.} The no-privacy reference (lambdas
$\lambda_1=\lambda_3=\lambda_4=0$, $\lambda_2=1$: an encoder trained for
next-location prediction alone, with no reconstruction or
re-identification adversary) measures the disparity already present in
an undefended learned representation, decoupled from any privacy
objective. On \textbf{Foursquare}, the residual-risk ratio rises
clearly and significantly from the no-privacy baseline to plain Mo-PAE
($3.40\times$, $[2.74,4.32]$, to $15.54\times$, $[10.84,23.71]$ --
non-overlapping CIs), while the absolute difference (decile 0 $-$
decile 9) moves the \emph{opposite} direction, from $0.505$,
$[0.430,0.577]$, to $0.456$, $[0.396,0.519]$ -- overlapping CIs, not
independently significant, but notably not widening either. Decile 0's
own post-mechanism risk drops substantially in absolute terms
($0.715 \to 0.487$). On \textbf{Gowalla-NYC}, neither the ratio
($3.99\times,[2.63,6.27] \to 6.93\times,[4.24,11.71]$) nor the
difference ($0.446,[0.294,0.597] \to 0.484,[0.336,0.636]$) is
statistically distinguishable between the no-privacy baseline and
plain Mo-PAE -- both metrics' confidence intervals overlap
substantially, so ``widens'' is not a safe claim here despite the point
estimates moving upward. On \textbf{GeoLife}, both metrics agree the
two are indistinguishable. The one claim that survives across datasets
is narrower than a single aggregate statistic would suggest:
Foursquare's residual-risk ratio widens clearly and significantly from
the undefended baseline to plain Mo-PAE; nothing else in this
comparison does, by either metric, on any dataset.

\textbf{A fairness term built on the mechanism's own co-trained
adversary targets the wrong decile.} We test this directly with
floor\_loss, a fairness term that penalizes the co-trained
re-identification unit's confidence once it exceeds a threshold. A
threshold set in absolute terms is unreachable if the unit's own
confidence never gets that high: on GeoLife, that confidence stays
within $0.005$--$0.008$ throughout training, regardless of decile,
because the encoder's own adversarial training suppresses it near
chance by construction -- so a fixed, absolute threshold gives the term
exactly zero gradient at every step. Replacing the fixed threshold with
one scaled to the chance level restores a genuinely nonzero gradient.
Table~\ref{tab:04d-decile-reduction}'s floor\_loss rows show what that
restored gradient actually does: decile 0's relative reduction stays
statistically unchanged from plain Mo-PAE at every fairness weight
tested, while decile 9's reduction rises monotonically with the
fairness weight and the residual-risk ratio gets \emph{worse}, not
better. The restored gradient targets whichever user's in-model
confidence crosses the (now chance-relative) threshold first, and
Table~\ref{tab:inmodel-vs-external} shows that is not preferentially
decile 0's users: GeoLife's in-model adversary stays within a narrow
band across all ten deciles (far flatter than the external attacker's
much wider range), so a confidence-threshold term built on it cannot
reliably single out the decile this paper's own evaluation protocol
identifies as worst-off. This is the mechanistic account, not just the
empirical one, of why a loss-based individual-fairness term built on an
in-model adversary struggles to target real re-identification risk.

A leveling-down check (testing, per decile, whether users harmed
relative to plain Mo-PAE at matched utility are disproportionately the
ones already best-protected, with harmed/helped thresholds derived from
a seed-to-seed null rather than a fixed epsilon) found no comparison
in this batch where the harmed rate clears that null -- most likely an
underpowered test given only 3 encoder seeds, not evidence that no such
redistribution is happening; Section~\ref{sec:phase4} reports a
better-powered version of this check (10 encoder seeds) that does find
a real effect at the aggregate (decile-0) level.

\textbf{A ceiling reference is only informative if it is itself trained
to convergence.} Every mechanism in this section uses a fixed training
budget with no early stopping; a per-epoch diagnostic on GeoLife found
all three mechanisms (no-privacy, plain, and the fairness-weighted
variant tested in Section~\ref{sec:phase4}) still rising substantially
partway through that budget, with their relative utility order changing
by the time training is extended. Table~\ref{tab:no-privacy-ceiling}
reports the no-privacy ceiling both at the original budget and once
extended to convergence: the ceiling rises on GeoLife (0.623 $\to$
0.742) but \emph{falls} on Foursquare (0.280 $\to$ 0.214) and
Gowalla-NYC (0.322 $\to$ 0.251) -- both corrected values fit smoothly
into their own dataset's privacy-weight sweep (Section~\ref{sec:phase4}),
so the shorter-budget values were inflated by the no-privacy
objective's faster early convergence, not a regression in the corrected
values.

\begin{table}[t]
\centering
\caption{No-privacy ceiling utility, before (20 epochs, 04d) and after (40 epochs, Checkpoint 04e) correcting for undertraining. Generated from \texttt{reports/figures/phase4e\_analysis/reference\_curve\_comparisons.json}.}
\label{tab:no-privacy-ceiling}
\small
\begin{tabular}{@{}lcc@{}}
\toprule
Dataset & Ceiling @ 20 epochs (stale) & Ceiling @ 40 epochs (corrected) \\
\midrule
GeoLife & 0.623$\pm$0.014 & 0.742$\pm$0.005 \\
Foursquare & 0.280$\pm$0.002 & 0.214$\pm$0.002 \\
Gowalla-NYC & 0.322$\pm$0.005 & 0.251$\pm$0.001 \\
\bottomrule
\end{tabular}
\end{table}

